%% kapa: exam=1 round=29 year=2018 period=2 form=A subject=accounting subject_ko=회계학 seq=34 number=74
%% source: Q-Net 감정평가사 기출문제 (cst00309) artlSeq=5202808 file=extracted/1008/5202808/2018년 제29회 감정평가사 1차시험 기출문제/2018년 제29회 감정평가사 1차시험 기출문제/2교시 A형.pdf
%% conversion: kordoc 4.24.51 -> tools/gen1.py (auto); review=proofread-sonnet; flags=none
\begin{question}{74}
(주)감평은 수선부문과 동력부문의 두 개의 보조부문과 도색부문과 조립부문의 두 개의 제조부문으로 구성되어 있다. (주)감평은 상호배부법을 사용하여 보조부문의 원가를 제조부문에 배부한다. 20x1년도 보조부문의 용역제공은 다음과 같다.
\par{\small\centering
\begin{tabular}{|c|c|c|c|c|}
\hline
\multirow{2}{*}{제공부문} & \multicolumn{2}{c|}{보조부문} & \multicolumn{2}{c|}{제조부문} \\ \cline{2-5}
 & 수선 & 동력 & 도색 & 조립 \\ \hline
수선(시간) & - & 400 & 1,000 & 600 \\ \hline
동력(kwh) & 2,000 & - & 4,000 & 4,000 \\ \hline
\end{tabular}\par}
20x1년도 보조부문인 수선부문과 동력부문으로부터 도색부문에 배부된 금액은 ￦100,000이고, 조립부문에 배부된 금액은 ￦80,000이었다. 동력부문의 배부 전 원가는? \par
\begin{choices}
  \item ￦75,000
  \item ￦80,000
  \item ￦100,000
  \item ￦105,000
  \item ￦125,000
\end{choices}
\end{question}
\answer{4}
